\documentclass[10pt,a4paper]{amsart}
\usepackage[margin=19mm]{geometry}
\usepackage{amsmath,amssymb,mathtools}
\usepackage[scr=rsfs,cal=euler]{mathalfa}
\usepackage{xfrac}
\usepackage[unicode,hidelinks,bookmarks=false]{hyperref}
\newcommand{\ie}{i.e.\ }
\newcommand{\cf}{cf.\ }
\newcommand{\resp}{respectively\ }
\newcommand{\Subsection}{\subsection}
\providecommand{\nobreakdash}{\nobreak\hbox{-}}
\setlength{\emergencystretch}{2em}
\multlinegap0pt
\usepackage{amssymb}
\usepackage{mathtools}
\usepackage[shortlabels]{enumitem}
\usepackage{float}
\usepackage{tikz}
\usepackage{xcolor}
\newtheorem{lemma}{Lemma}
\usepackage{cleveref}
\crefname{lemma}{Lemma}{Lemmas}
\newcommand{\Ber}{\textup{Ber}}
\newcommand{\Coef}{\textup{Coef}}
\newcommand{\E}{\mathbb{E}}
\newcommand{\EE}{\mathbb E}
\newcommand{\F}{\mathbb{F}}
\newcommand{\FF}{\mathbb F}
\newcommand{\abs}[1]{\left|#1\right|}
\def\eqdef{\overset{\rm{def}}{=}}
\renewcommand{\ge}{\geqslant}
\renewcommand{\geq}{\ge}
\renewcommand{\le}{\leqslant}
\renewcommand{\leq}{\le}
\newcommand{\rank}{\textup{rank}}
\title{Vertex-minor universality of a random graph}
\author{Ting-Wei Chao, Zixuan Xu}
\date{Source: arXiv:2603.13600v3 (CC BY 4.0)}
\begin{document}
\maketitle

\noindent\emph{Source and licence.} Vertex-minor universality of a random graph. Version arXiv:2603.13600v3 (CC BY 4.0), distributed by arXiv under the Creative Commons Attribution 4.0 International licence, http://creativecommons.org/licenses/by/4.0/, as recorded in the arXiv licence metadata for this exact version and shown on the abstract page https://arxiv.org/abs/2603.13600v3. Full licence notice: \texttt{licenses/arxiv-2603.13600-CC-BY-4.0.txt}.\par\smallskip

\noindent\textit{Editorial scope.} Counted unit: the article's lemma lem:1 (source0.tex lines 275--278). The article's complete original proof of it is delivered with it (source0.tex lines 280--335, one \texttt{proof} environment of 56 lines). No mathematics is written, repaired or re-derived: the statement and the proof are the article's own lines, in the article's own notation, and no retained line is substituted or re-flowed.  No further source-local context is needed: every cross-reference of every retained line resolves inside the two delivered spans.  Nothing was omitted from any retained span, so the package carries no omission record.  No retained line contains a citation, so the wrapper carries no bibliography and no bibliography entry.  No retained line contains a figure environment, an image-inclusion command or drawing code, so no image is delivered and the package declares no asset dependency.  Editorial formatting only, in addition: an AMS \texttt{amsart} preamble that restates the article's own declarations and macros used by the retained lines, namely the environments \texttt{lemma} on separate counters, together with the standard packages those lines need.  The retained environments print in this document as Lemma 1 in order of appearance, whereas the article prints the same items with the numbering of its own full document.  The complete native source of the article as distributed in the arXiv e-print for this exact version is bundled unchanged as \texttt{sources/196344/source0.tex}.
\begin{lemma}\label{lem:1}
    Let $p\in [0,1]$ and $q\eqdef\min(p,1-p)$. Let $r$ be a positive integer and $X\in \F_2^m$ be a vector with independent $\Ber(p)$ entries. Let $f\in\FF_2[x_1,\dots,x_m]$ be a multilinear polynomial with $\deg f\leq 2$. Suppose that $\rank(\Coef_2(f))=r$. Then we have
    \[\abs{\E[(-1)^{f(X)}]}\le (1-2q^2)^{\lfloor r/6\rfloor}.\]
\end{lemma}

\begin{proof}
    We prove by induction on $r$. In the base case where $r\leq 6$, the statement is trivial since the right hand side is $1$.

    Now suppose the statement is true for all integers less than $r$ and $r\geq 6$. Let $X=(X_1,\dots,X_m)\in \F_2^m$ be a vector with independent $\Ber(p)$ entries.
    Let $f\in\FF_2[x_1,\dots,x_m]$ be a multilinear polynomial with $\deg f\leq 2$ such that $\rank(\Coef_2(f))=r$. Let $\Coef_2(f)=(a_{ij})_{i,j=1}^m$ and we can further assume without loss of generality that $a_{12}=1$. Then we can write $f$ as
    \[f=x_1x_2+x_1f_1+x_2f_2+f_3,\]
    where $f_1,f_2,f_3\in \FF_2[x_3,\dots,x_m]$ such that $\deg f_1,\deg f_2\leq 1$, $\deg f_3\leq 2$, and $f_3$ is multilinear. Note that $\Coef_2(f_3)$ is a minor of $\Coef_2(f)$ obtained from removing the first two rows and columns from $\Coef_2(f)$. Therefore, we have $\rank(\Coef_2(f_3))\geq r-4$. Denote $X'=(X_3,\dots,X_m)$ and we have 
    \[\EE[(-1)^{f(X)}] = \EE_{X'}\left[(-1)^{f_3(X')}\EE_{X_1,X_2}[(-1)^{X_1X_2+X_1f_1(X')+X_2f_2(X')}\mid X']\right].\]

    We first prove the case $q=p\leq 1/2$.
    By a direct calculation, for any $a,b\in\FF_2$ and since $X_1,X_2\sim \Ber(p)$ are independent, we have
    \begin{align}\label{eq:Exp}
        \EE_{X_1,X_2}[(-1)^{X_1X_2+aX_1+bX_2}] = (1-p)^2 + p(1-p)((-1)^a+(-1)^b)-p^2(-1)^{a+b}.
    \end{align}
    Using the identity that 
    \[(-1)^{a+b}= 1+((-1)^a + (-1)^b)-2(-1)^{ab},\]
    we obtain
    \begin{equation}
        \EE_{X_1,X_2}[(-1)^{X_1X_2+aX_1+bX_2}] = (1-2p) + p(1-2p)((-1)^a+(-1)^b)+2p^2(-1)^{ab}.
    \end{equation}
    Therefore, we can write
    \begin{align*}
        \EE[(-1)^{f(X)}] =&\,\EE_{X'}\left[(-1)^{f_3(X')}\left((1-2p) + p(1-2p)((-1)^{f_1(X')}+(-1)^{f_2(X')})+2p^2(-1)^{f_1(X')f_2(X')}\right)\right]\\
        =&\,(1-2p)\EE_{X'}\left[(-1)^{f_3(X')}\right]+p(1-2p)\EE_{X'}\left[(-1)^{f_3(X')+f_1(X')}\right]\\
        &+p(1-2p)\EE_{X'}\left[(-1)^{f_3(X')+f_2(X')}\right]+2p^2\EE_{X'}\left[(-1)^{f_3(X')+f_1(X')f_2(X')}\right].
    \end{align*}
    Now, we can apply the inductive hypothesis to bound each term. Recall that we have $\rank(\Coef_2(f_3))\geq r-4$. Therefore, we obtain 
    \[\abs{\EE_{X'}\left[(-1)^{f_3(X')}\right]}\leq (1-2p^2)^{\lfloor(r-4)/6\rfloor}\leq (1-2p^2)^{\lfloor r/6\rfloor-1}.\]
    Furthermore, since $\deg f_1, \deg f_2\le 1$, we have $\Coef_2(f_3+f_1)=\Coef_2(f_3+f_2)=\Coef_2(f_3)$, which has rank at least $r-4$. Thus, we can conclude similarly that
    \[\abs{\EE_{X'}\left[(-1)^{f_3(X')+f_1(X')}\right]}\leq (1-2p^2)^{\lfloor r/6\rfloor-1}\text{ and }\abs{\EE_{X'}\left[(-1)^{f_3(X')+f_2(X')}\right]}\leq (1-2p^2)^{\lfloor r/6\rfloor-1}.\]
    Finally, note that the homogeneous degree $2$ part of $f_1f_2$ is $\sum_{i,j=3}^m\alpha_i\beta_jx_ix_j$ if $f_1=\sum_{i=3}^m\alpha_ix_i+c_1$, and $f_2=\sum_{j=3}^m\beta_jx_j+c_2$. Consider the multilinearization $g$ of $f_1f_2$, i.e. the polynomial $g$ obtained by replacing all the terms $x_i^2$ by $x_i$ in $f_1f_2$. It follows that the coefficient of $x_ix_j$ in $g$ is $\alpha_i\beta_j+\alpha_j\beta_i$ if $i\neq j$. Therefore, we have
    \[\Coef_2(g)=\alpha^T\beta+\beta^T\alpha\]
    where $\alpha=(\alpha_3,\dots,\alpha_m)$ and $\beta=(\beta_3,\dots,\beta_m)$. Since the diagonal terms of $\alpha^T\beta+\beta^T\alpha$ are always $0$, they agree with the diagonal terms of $\Coef_2(g)$.
    Therefore, we have $\rank(\Coef_2(g))\leq 2$. Hence, we obtain that $\rank(\Coef_2(f_3+g))\geq r-6$. By the inductive hypothesis, we have
    \[\abs{\EE_{X'}\left[(-1)^{f_3(X')+f_1(X')f_2(X')}\right]}=\abs{\EE_{X'}\left[(-1)^{f_3(X')+g(X')}\right]}\leq (1-2p^2)^{\lfloor(r-6)/6\rfloor} = (1-2p^2)^{\lfloor r/6\rfloor-1}.\]
    Combining the above inequalities and the fact that $p,(1-2p)\geq 0$ for $p\le 1/2$, we obtain
    \[\abs{\EE[(-1)^{f(X)}]}\leq \left((1-2p)+p(1-2p)+p(1-2p)+2p^2\right)(1-2p^2)^{\lfloor r/6\rfloor-1}=(1-2p^2)^{\lfloor r/6\rfloor}.\]

    For the case $p\geq 1/2$ and $q=1-p$, we can combine \cref{eq:Exp} with the identity
    \[1=(-1)^{a+b}-((-1)^a + (-1)^b)+2(-1)^{ab}\]
    and conclude that
    \begin{align*}
        \EE_{X_1,X_2}[(-1)^{X_1X_2+aX_1+bX_2}] =&-(2p-1)(-1)^{a+b}+ (1-p)(2p-1)((-1)^a+(-1)^b)+2(1-p)^2(-1)^{ab}\\
        =& -(1-2q)(-1)^{a+b}+q(1-2q)((-1)^a+(-1)^b)+2q^2(-1)^{ab}.      
    \end{align*}
    Similar as before, we have
    \begin{align*}
        &\EE[(-1)^{f(X)}]\\
        &=\EE_{X'}\left[(-1)^{f_3(X')}\left(-(1-2q)(-1)^{f_1(X')+f_2(X')} + q(1-2q)((-1)^{f_1(X')}+(-1)^{f_2(X')})+2q^2(-1)^{f_1(X')f_2(X')}\right)\right]\\
        &=-(1-2q)\EE_{X'}\left[(-1)^{f_3(X')+f_1(X')+f_2(X')}\right]+q(1-2q)\EE_{X'}\left[(-1)^{f_3(X')+f_1(X')}\right]\\
        &\quad+q(1-2q)\EE_{X'}\left[(-1)^{f_3(X')+f_2(X')}\right]+2q^2\EE_{X'}\left[(-1)^{f_3(X')+f_1(X')f_2(X')}\right].\\
        &\leq \left((1-2q)+q(1-2q)+q(1-2q)+2q^2\right)(1-2q^2)^{\lfloor r/6\rfloor-1}=(1-2q^2)^{\lfloor r/6\rfloor},
    \end{align*}
    where the last inequality follows from applying the inductive hypothesis to $f_3+f_1+f_2$, $f_3+f_1$, $f_3+f_2$, and $f_3+g$, where $g$ is the multilinearization of $f_1f_2$. We also used the fact that $q,(1-2q)\geq 0$.
    This concludes the proof.
\end{proof}

\end{document}
